% !PRINTED 172
\weberpage{172}

Wenn wir also in (17) $\omega$ durch $2\omega$ ersetzen und die Quadratwurzel ziehen, so folgt:

\weq{19}{q = \frac{1}{4}\frac{\varkappa_1}{\varkappa_1'}\, e^{-\frac{1}{2}\frac{G(\varkappa_1{}^2)}{F(\varkappa_1{}^2)}}.}

Hierin erstreckt sich nun der Konvergenzbereich \"uber den in der Ebene $\varkappa_1{}^2$ gelegenen Einheitskreis und dieser entspricht der ganzen Ebene $\varkappa^2$. Denn der reelle Teil vom $\varkappa'$ ist in der \gesperrt{ganzen Ebene} der $\varkappa^2$ positiv mit Ausnahme der beiden Ufer des von $+1$ nach $+\infty$ verlaufenden Schnittes, an denen $\varkappa'$ rein imagin\"ar ist. Daraus ergibt sich, da\ss{} der absolute Wert von $\varkappa_1$ kleiner als 1 ist, und da\ss{} die Peripherie des Einheitskreises in der Ebene $\varkappa_1$ den beiden Ufern dieses Schnittes entspricht.

Man kann aber die Konvergenz dieser Ausdr\"ucke noch verbessern durch eine abermalige Anwendung der \gesperrt{Landenschen} Transformation. Bezeichnen wir das Resultat einer nochmaligen Verdoppelung von $\omega$ durch den Index 2, so folgt:

\weq{20}{\begin{aligned}
\sqrt{\varkappa_2} &= \frac{1-\sqrt{\varkappa'}}{1+\sqrt{\varkappa'}}, \\
4K_2 &= (1+\sqrt{\varkappa'})^2 K, \quad K_2' = (1+\sqrt{\varkappa'})^2 K', \\
q &= \sqrt{\frac{\varkappa_2}{4\varkappa_2'}}\, e^{-\frac{1}{4}\frac{G(\varkappa_2{}^2)}{F(\varkappa_2{}^2)}}.
\end{aligned}}

worin nun der Konvergenzbereich die Ebene der $\varkappa^2$ \gesperrt{zweimal} (mit einer Verzweigung im Punkte $\varkappa^2=1$) bedeckt.

Diese Operation kann man fortsetzen, indem man nach der Formel (17) aus $\varkappa_2$ durch Verdoppelung von $\omega$ eine neue Gr\"o\ss{}e $\varkappa_3$, daraus ebenso $\varkappa_4$ usf. herleitet. Man bekommt dann eine Reihe von Ausdr\"ucken f\"ur $q$, deren Konvergenz eine immer bessere wird. F\"ur ein beliebiges $\nu$ ist:

\weq{21}{q = \sqrt[2^{\nu-1}]{\frac{\varkappa_\nu}{4\varkappa_\nu'}}\, e^{-\frac{G(\varkappa_\nu{}^2)}{2^\nu F(\varkappa_\nu{}^2)}}.}

Die erste Gleichung (18)

\[
\varkappa_1 = \frac{\varkappa^2}{(1+\varkappa')^2}
\]

zeigt, da\ss{}, so lange der absolute Wert von $\varkappa^2$ kleiner als 1 ist, und folglich der absolute Wert von $1+\varkappa'$ gr\"o\ss{}er als 1, der absolute Wert von $\varkappa_1$ kleiner ist als der von $\varkappa^2$, und folglich
%CONT>>
